\documentclass[sigconf,nonacm]{acmart}

%%% Remove ACM-specific elements for proposal submission
\settopmatter{printacmref=false}
\renewcommand\footnotetextcopyrightpermission[1]{}
\pagestyle{plain}

\begin{document}

\title{From Retrieval to Generation: Comparing Explainable, Retrieval-Augmented, and LLM-Native Text Recommender Systems}
\subtitle{Tutorial Proposal for RecSys 2026}

\author{Yejin Seo}
\email{yseo@legalzoom.com}
\affiliation{%
  \institution{LegalZoom; Harvard University}
}

\author{Sarthak Baral}
\email{sab7803@g.harvard.edu}
\affiliation{%
  \institution{Harvard University}
}

\author{Kaylee Vo}
\email{kvo@seas.harvard.edu}
\affiliation{%
  \institution{Harvard University}
}

\begin{abstract}
Recent progress in transformer encoders, retrieval-augmented generation, and large language models has created multiple viable ways to build text-rich recommender systems.
Yet the field still lacks a practitioner-oriented tutorial that compares these choices under a common framework.
This tutorial presents a structured, empirical, and implementation-aware comparison of three paradigms for text-driven recommendation:
collaborative-evidence explainable recommenders, retrieval-augmented recommenders with LLM prompting or reranking, and LLM-native generative recommenders.

Rather than treating these paradigms in isolation, the tutorial shows how they differ in recommendation quality, explanation quality, operational complexity, latency, and cost.
We discuss how to construct comparable pipelines, how to evaluate ranking performance with standard top-$K$ metrics, and how to evaluate explanations for relevance, specificity, and hallucination risk.
We also present a practical ``when to use what'' framework that helps attendees map method choice to product constraints such as sparse interactions, cold start, explainability requirements, and serving budgets.

The tutorial is aimed at researchers and practitioners who want a grounded path from modern language models to deployable recommendation systems.
Attendees will leave with conceptual clarity, implementation guidance, and a reusable evaluation template for future work.
\end{abstract}

\maketitle

\section{Tutorial Length}
Half-day tutorial delivered over two conference tutorial blocks, for approximately 3 hours of instructional time plus the conference coffee break.

\section{Motivation}
Language models are rapidly changing the design space of recommender systems, especially in domains where item descriptions, reviews, metadata, and other textual signals are abundant.
However, the practical choice among model families remains unclear.
Teams can now build:
\begin{enumerate}
    \item recommenders that stay tightly grounded in collaborative evidence and generate interpretable explanations from historical behavior,
    \item retrieval-augmented systems that combine semantic retrieval with LLM reasoning or reranking, or
    \item fully generative recommenders that use instruction-following and language generation as the main recommendation interface.
\end{enumerate}

These options differ not only in offline ranking performance, but also in explanation quality, controllability, serving cost, and deployment complexity.
Existing overviews often focus on a single paradigm at a time.
This tutorial instead offers a controlled comparative perspective, designed to help the RecSys community reason about trade-offs rather than chase a single fashionable architecture.

The tutorial is timely for RecSys because it connects several active subareas that the community is already engaging with---explainability, retrieval, multimodal/text-rich recommendation, and LLM-based recommendation---while keeping the emphasis on practical method choice and reproducible evaluation.

\section{Instructors}
\textbf{Yejin Seo} is a Principal Data Scientist at LegalZoom and a student at Harvard University.
\textbf{Sarthak Baral} is a student at Harvard University.
\textbf{Kaylee Vo} is a Graduate Researcher and student at Harvard University.

\section{Detailed Outline}

\subsection{Block A (90 minutes)}

\subsubsection{Framing the Problem Space}
\begin{itemize}
    \item Why text-rich recommendation matters now
    \item Where collaborative filtering stops being enough
    \item Why explanation quality matters in product settings
    \item The three paradigms covered in this tutorial
    \item What ``controlled comparison'' means in practice
\end{itemize}

\subsubsection{Representations and Retrieval Foundations}
\begin{itemize}
    \item Item and user text representations
    \item Transformer encoders for semantic retrieval
    \item Pretrained embeddings versus task-adapted embeddings
    \item Candidate generation, reranking, and grounding pathways
    \item Failure modes in sparse and noisy text settings
\end{itemize}

\subsubsection{Paradigm 1: Explainable Recommenders Grounded in Collaborative Evidence}
\begin{itemize}
    \item Architecture and data flow
    \item How collaborative signals and textual signals are combined
    \item Explanation construction from historical user--item evidence
    \item Template-based versus neural explanation generation
    \item \textit{Advantages:} stronger grounding, easier justification, tighter connection to user history
    \item \textit{Limitations:} weaker cold-start behavior, dependence on interaction density, narrower generalization outside known catalogs
    \item Example evaluation setup and failure analysis
\end{itemize}

\subsubsection{Paradigm 2: Retrieval-Augmented Recommenders}
\begin{itemize}
    \item Dense retrieval for initial candidate generation
    \item Prompt assembly and context injection
    \item LLM reranking versus LLM explanation generation
    \item Retrieval quality versus generation quality trade-offs
    \item Controlling hallucinations through grounding
    \item Serving considerations: batching, caching, latency budgets, cost-aware inference
    \item Design decisions that most affect production viability
\end{itemize}

\subsection{Block B (90 minutes)}

\subsubsection{Paradigm 3: LLM-Native Generative Recommenders}
\begin{itemize}
    \item End-to-end recommendation through generation
    \item Prompt-based user modeling
    \item Few-shot and zero-shot recommendation patterns
    \item \textit{Strengths:} flexibility, low engineering overhead at prototype time, strong natural-language interaction
    \item \textit{Limitations:} output instability, grounding problems, ranking/generation mismatch, high serving cost
\end{itemize}

\subsubsection{Unified Experimental Framework}
\begin{itemize}
    \item Shared candidate pools and fair comparisons
    \item Recommended benchmark structure using text-rich recommendation data
    \item Ranking metrics: HR@10, NDCG@10, Recall@$K$
    \item Explanation metrics: \emph{relevance} (does the explanation cite items or attributes the user actually interacted with?), \emph{specificity} (does it go beyond generic praise to reference concrete item features?), \emph{consistency} (does the same input produce semantically stable explanations across runs?), and \emph{hallucination rate} (does the explanation assert facts not supported by the input data?)
    \item System metrics: latency, token cost, implementation complexity, reproducibility burden
    \item How to interpret contradictory metric outcomes
\end{itemize}

\subsubsection{Hands-On Walkthrough}
\begin{itemize}
    \item Minimal embedding-based retrieval pipeline
    \item Minimal retrieval-augmented recommendation pipeline
    \item Example explanation evaluation
    \item Notebook structure for reproducibility
    \item How to extend the starter code after the tutorial
\end{itemize}

\subsubsection{Decision Framework and Open Challenges}
\begin{itemize}
    \item Which approach to use under: sparse interactions, cold-start items, explanation-heavy product requirements, strict serving budgets, fast prototyping constraints
    \item Bias, safety, and evaluation gaps
    \item Research opportunities in grounded explanation and benchmark design
    \item Closing Q\&A
\end{itemize}

\section{Target Audience and Prerequisites}
This tutorial is designed for an intermediate to advanced audience.
It will be most useful for recommender-systems researchers, PhD and advanced Master's students, ML engineers and applied scientists working with ranking, retrieval, or LLM systems, and practitioners interested in explainability and text-rich recommendation settings.

\paragraph{Expected background.}
Familiarity with standard recommender-system concepts; comfort with basic ranking metrics; basic machine-learning knowledge.
Some exposure to transformers or embeddings is helpful but not required.

\section{Materials and Logistics}

\paragraph{Materials for attendees.}
Slide deck in PDF form; starter notebooks for the main pipelines; code templates for evaluation; reading list; reproducibility guide; post-tutorial repository containing notebooks and sample prompts.

\paragraph{Attendee needs.}
A laptop is recommended for the optional hands-on component.
No specialized hardware is required for following the tutorial.

\paragraph{Instructor needs.}
Standard audio-video setup and internet access.

\section{Teaching Experience}
\textbf{Yejin Seo:} Internal ML training curriculum design and delivery at LegalZoom; hands-on model education for cross-functional and data-science audiences.

\textbf{Sarthak Baral:} Taught undergraduate-level courses in Software Development Lifecycle at King's College London.

\textbf{Kaylee Vo:} Teaching Fellow for Harvard College, Applied Mathematics~115 (Mathematical Modeling), an upper-level undergraduate course.
Developed lecture notes\footnote{\url{https://drive.google.com/file/d/1NofKRNLueNZELv7cz7GXYgYVtnXb4cC5/view}} and code materials.\footnote{\url{https://github.com/zhimingkuang/Harvard-AM-115/blob/main/11_diffusion/solve_diffusion_with_pinn_bfgs.ipynb}}
Also taught introductory economics at Stanford Pre-Collegiate Summer Institutes.

\section{Relevant Publications}
The instructors are actively conducting graduate research at Harvard University on text-rich recommendation, explainable recommendation, and LLM-based ranking systems, with multiple manuscripts currently under review at peer-reviewed venues.
Seo additionally brings applied experience designing and deploying recommendation and ranking systems in production at LegalZoom.
The tutorial content is grounded in a comparative implementation study that the team is developing as part of this research program; the codebase and evaluation framework will be released alongside the tutorial materials.

\section{Differences from Prior Tutorials}

Compared with \textit{Agentic LLM for Recommender Systems} (RecSys~2025), this tutorial is not centered on agentic behavior or modular agent design; its core contribution is a comparative framework for choosing among grounded explainable recommenders, retrieval-augmented recommenders, and generative recommenders.

Compared with \textit{Multi Agentic Recommender Systems} (RecSys~2025), our focus is on controlled evaluation, explanation quality, and deployment trade-offs across model families rather than multi-agent orchestration patterns.

Compared with \textit{Standard Practices for Data Processing and Multimodal Feature Extraction in Recommendation with DataRec and Ducho} (RecSys~2025), this tutorial concentrates on model-family selection and end-to-end decision-making in text-rich recommendation, not preprocessing pipelines or multimodal feature engineering.

Compared with \textit{Recent Advances in Generative Conversational Recommender Systems} (RecSys~2025), this tutorial is broader than conversation and explicitly includes non-generative, evidence-grounded baselines in the comparison.

Compared with \textit{A Tutorial on Feature Interpretation in Recommender Systems} (RecSys~2024), this tutorial focuses on choosing among three modern architecture families and evaluating both recommendation outputs and generated explanations, rather than feature interpretation within existing pipelines.

\bibliographystyle{ACM-Reference-Format}

\begin{thebibliography}{7}

\bibitem{harper2015movielens}
F.~Maxwell Harper and Joseph~A. Konstan. 2015.
\newblock The MovieLens Datasets: History and Context.
\newblock {\em ACM Transactions on Interactive Intelligent Systems\/} 5, 4, Article~19 (2015), 19~pages.

\bibitem{reimers2019sentence}
Nils Reimers and Iryna Gurevych. 2019.
\newblock Sentence-BERT: Sentence Embeddings using Siamese BERT-Networks.
\newblock In {\em Proceedings of the 2019 Conference on Empirical Methods in Natural Language Processing (EMNLP)}.

\bibitem{lewis2020rag}
Patrick Lewis, Ethan Perez, Aleksandra Piktus, Fabio Petroni, Vladimir Karpukhin, Naman Goyal, Heinrich K\"{u}ttler, Mike Lewis, Wen-tau Yih, Tim Rockt\"{a}schel, Sebastian Riedel, and Douwe Kiela. 2020.
\newblock Retrieval-Augmented Generation for Knowledge-Intensive NLP Tasks.
\newblock In {\em Advances in Neural Information Processing Systems\/} 33.

\bibitem{zhang2020explainable}
Yongfeng Zhang and Xu~Chen. 2020.
\newblock Explainable Recommendation: A Survey and New Perspectives.
\newblock {\em Foundations and Trends in Information Retrieval\/} 14, 1 (2020), 1--101.

\bibitem{zhang2020bertscore}
Tianyi Zhang, Varsha Kishore, Felix Wu, Kilian~Q. Weinberger, and Yoav Artzi. 2020.
\newblock BERTScore: Evaluating Text Generation with BERT.
\newblock In {\em Proceedings of the Eighth International Conference on Learning Representations (ICLR)}.

\bibitem{shehmir2025llm4rec}
Syed~Ali Shehmir and Rasha Kashef. 2025.
\newblock LLM4Rec: A Comprehensive Survey on the Integration of Large Language Models in Recommender Systems---Approaches, Applications and Challenges.
\newblock {\em Information Fusion\/} 117 (2025), 102842.

\end{thebibliography}

\end{document}
